% GENERATED FILE. Edit canonical lessons, then run npm run generate:books.
% canonical-source-hash: fnv1a64:54415c0c0b3b2c0d
% canonical-lessons: ZH-C125-shengqi, ZH-C125-danxin, ZH-C125-fangxin, ZH-C125-zhaoji, ZH-C125-nanguo

\chapter{Angry and Worried}
\label{ch:zh-angry-and-worried}
\hlchaptermodality{125}

\begin{chapteropening}
\textbf{By the end of this chapter:} \emph{I can give a feeling as a reason: angry, worried, at ease, anxious or upset.}

Complete the last lesson of chapter 125: I can give a feeling as a reason: angry, worried, at ease, anxious or upset.
\end{chapteropening}

\section[\texorpdfstring{\zh{生气} (shēngqì)}{生气 (shēngqì)}]{\zh{生气} (shēngqì) — to get angry}

\label{lesson:ZH-C125-shengqi}



\begin{quote}
Before the new one: say the Mandarin for ugly, then the Mandarin for to resemble.
\end{quote}

\subsection*{You'll want to know: \zh{生气}}
\textbf{\zh{生气}} — \emph{shēngqì} — \textquotedblleft{}to get angry\textquotedblright{}.

\begin{quote}
\textbf{\zh{因为我没有来他生气了}} \emph{yīnwèi wǒ méiyǒu lái tā shēngqì le} — He got angry because I didn't come.
\end{quote}

\begin{cousinweb}
\textbf{\zh{生气}} (\emph{shēngqì}) is built from \zh{生} 'to grow' and \zh{气} 'breath', both inherited from Middle Chinese, \emph{srjaeng} and \emph{khj+jH} (Guangyun, in Baxter's transcription). The sense angry is not given there, so how the pair came to it is not settled.
\end{cousinweb}

\begin{grammarlens}[title={What you've built}]
One more word for giving reasons.
\end{grammarlens}

\subsection*{Your turn}
\begin{itemize}
\raggedright
  \item \emph{Say it:} \emph{shēngqì}
  \item \emph{Say it:} \emph{shēngqì}, once more
  \item \emph{Say it:} \emph{yīnwèi wǒ méiyǒu lái tā shēngqì le} — He got angry because I didn't come.
  \item \emph{Recall:} say the Mandarin for ugly, then the Mandarin for to resemble, then say \emph{shēngqì} again
\end{itemize}

\begin{tcolorbox}[breakable,colback=teal!4,colframe=teal!35!black,title={Before you move on}]
What does \zh{生气} mean? (\textquotedblleft{}To get angry\textquotedblright{}.) Say it once more.
\end{tcolorbox}
\section[dānxīn]{dānxīn — to worry}

\label{lesson:ZH-C125-danxin}



\begin{quote}
Before the new one: say the Mandarin for to resemble, then the Mandarin for to get angry.
\end{quote}

\subsection*{You'll want to know: dānxīn}
\emph{dānxīn} — \textquotedblleft{}to worry\textquotedblright{}.

\begin{quote}
\emph{nǐ bú yào dānxīn} — Don't you worry.
\end{quote}

Its characters come later; for now it is a word for the ear.

\begin{cousinweb}
\emph{dānxīn} is built from \emph{dān} 'to carry on the shoulder' and \emph{xīn} 'heart', both inherited from Middle Chinese, \emph{tam} and \emph{sim} (Guangyun, in Baxter's transcription). The pair is not given there; a load carried in the heart is the plain picture, but the step is not settled.
\end{cousinweb}

\begin{grammarlens}[title={What you've built}]
One more word for giving reasons.
\end{grammarlens}

\subsection*{Your turn}
\begin{itemize}
\raggedright
  \item \emph{Say it:} \emph{dānxīn}
  \item \emph{Say it:} \emph{dānxīn}, once more
  \item \emph{Say it:} \emph{nǐ bú yào dānxīn} — Don't you worry.
  \item \emph{Recall:} say the Mandarin for to resemble, then the Mandarin for to get angry, then say \emph{dānxīn} again
\end{itemize}

\begin{tcolorbox}[breakable,colback=teal!4,colframe=teal!35!black,title={Before you move on}]
What does dānxīn mean? (\textquotedblleft{}To worry\textquotedblright{}.) Say it once more.
\end{tcolorbox}
\section[fàngxīn]{fàngxīn — to stop worrying; to be at ease}

\label{lesson:ZH-C125-fangxin}



\begin{quote}
Before the new one: say the Mandarin for to get angry, then the Mandarin for to worry.
\end{quote}

\subsection*{You'll want to know: fàngxīn}
\emph{fàngxīn} — \textquotedblleft{}to stop worrying; to be at ease\textquotedblright{}.

\begin{quote}
\emph{nǐ fàngxīn wǒ hěn hǎo} — Don't worry, I'm fine.
\end{quote}

Its characters come later; for now it is a word for the ear.

\begin{cousinweb}
\emph{fàngxīn} is built from \emph{fàng}, the word for putting things down, and \emph{xīn} 'heart', both inherited from Middle Chinese, \emph{pjangH} and \emph{sim} (Guangyun, in Baxter's transcription); the dictionary glosses the first 'to drive out' and 'to send away'.
\end{cousinweb}

\begin{grammarlens}[title={What you've built}]
One more word for giving reasons.
\end{grammarlens}

\subsection*{Your turn}
\begin{itemize}
\raggedright
  \item \emph{Say it:} \emph{fàngxīn}
  \item \emph{Say it:} \emph{fàngxīn}, once more
  \item \emph{Say it:} \emph{nǐ fàngxīn wǒ hěn hǎo} — Don't worry, I'm fine.
  \item \emph{Recall:} say the Mandarin for to get angry, then the Mandarin for to worry, then say \emph{fàngxīn} again
\end{itemize}

\begin{tcolorbox}[breakable,colback=teal!4,colframe=teal!35!black,title={Before you move on}]
What does fàngxīn mean? (\textquotedblleft{}To stop worrying; to be at ease\textquotedblright{}.) Say it once more.
\end{tcolorbox}
\section[zháojí]{zháojí — anxious; in a hurry}

\label{lesson:ZH-C125-zhaoji}



\begin{quote}
Before the new one: say the Mandarin for to worry, then the Mandarin for to be at ease.
\end{quote}

\subsection*{You'll want to know: zháojí}
\emph{zháojí} — \textquotedblleft{}anxious; in a hurry\textquotedblright{}.

\begin{quote}
\emph{wǒ hěn zháojí yīnwèi wǒ yǒushì} — I'm in a rush, because I've got something on.
\end{quote}

Its characters come later; for now it is a word for the ear.

\begin{cousinweb}
\emph{zháojí} is built from \emph{zháo} and \emph{jí}, both inherited from Middle Chinese, \emph{drjak} and \emph{kip} (Guangyun, in Baxter's transcription). The dictionary glosses the first 'to attach to' and the second 'urgent' and 'narrow'; how the pair came to mean anxious is not settled.
\end{cousinweb}

\begin{grammarlens}[title={What you've built}]
One more word for giving reasons.
\end{grammarlens}

\subsection*{Your turn}
\begin{itemize}
\raggedright
  \item \emph{Say it:} \emph{zháojí}
  \item \emph{Say it:} \emph{zháojí}, once more
  \item \emph{Say it:} \emph{wǒ hěn zháojí yīnwèi wǒ yǒushì} — I'm in a rush, because I've got something on.
  \item \emph{Recall:} say the Mandarin for to worry, then the Mandarin for to be at ease, then say \emph{zháojí} again
\end{itemize}

\begin{tcolorbox}[breakable,colback=teal!4,colframe=teal!35!black,title={Before you move on}]
What does zháojí mean? (\textquotedblleft{}Anxious; in a hurry\textquotedblright{}.) Say it once more.
\end{tcolorbox}
\section[nánguò]{nánguò — sad; upset}

\label{lesson:ZH-C125-nanguo}



\begin{quote}
Before the new one: say the Mandarin for to be at ease, then the Mandarin for anxious.
\end{quote}

\subsection*{You'll want to know: nánguò}
\emph{nánguò} — \textquotedblleft{}sad; upset\textquotedblright{}.

\begin{quote}
\emph{tā de gǒu bìng le tā hěn nánguò} — His dog is ill, and he's very upset.
\end{quote}

Its characters come later; for now it is a word for the ear.

\begin{cousinweb}
\emph{nánguò} is built from \emph{nán} 'hard' and \emph{guò} 'to pass', both inherited from Middle Chinese, \emph{nan} and \emph{kwaH} (Guangyun, in Baxter's transcription): hard to get through. The pair is not given there.
\end{cousinweb}

\begin{grammarlens}[title={What you've built}]
One more word for giving reasons.
\end{grammarlens}

\subsection*{Your turn}
\begin{itemize}
\raggedright
  \item \emph{Say it:} \emph{nánguò}
  \item \emph{Say it:} \emph{nánguò}, once more
  \item \emph{Say it:} \emph{tā de gǒu bìng le tā hěn nánguò} — His dog is ill, and he's very upset.
  \item \emph{Recall:} say the Mandarin for to be at ease, then the Mandarin for anxious, then say \emph{nánguò} again
\end{itemize}

\begin{tcolorbox}[breakable,colback=teal!4,colframe=teal!35!black,title={Before you move on}]
What does nánguò mean? (\textquotedblleft{}Sad; upset\textquotedblright{}.) Say it once more.
\end{tcolorbox}
